\section{From workflow submission to a scientifically acceptable artifact}
\label{sec:exec24-workflow-contracts}

A restarted causal batch needs an account of what was requested, what
actually ran and which result answers the banking question. A collections
team might request a monthly policy comparison, receive a trigger
acknowledgement, and later find either a completed simulation or no run at
all. A funding analyst might receive a Finished job whose process exited
unsuccessfully. Neither an acknowledgement nor a terminal state establishes
that the requested intervention contrast, population and information date
survived execution. The following historical source profiles supplement the
earlier client comparisons with immutable Dataiku client text, SAS REST
examples, a declared-range HTTP dependency profile and a concrete Posit
Local backend. They establish source contracts, not observed performance of
licensed installations.~\citep{exec24_workflow_contracts_dataiku,
exec24_workflow_contracts_sas,exec24_workflow_contracts_domino,
exec24_workflow_contracts_posit}

\subsection{Dataiku: the trigger denominator precedes the run denominator}

The authenticated Dataiku client identifies version 14.7.3. Its scenario
interface explicitly distinguishes a trigger firing from a scenario run.
A manual firing may be cancelled because a scenario is already running or
another trigger fires. The firing's run identifier is the trigger-run
identifier; it is not automatically the eventual scenario-run identifier.
The proposed adapter should retain both identities and their relation,
including the legitimate outcome that no scenario run is created. This
historical client cannot authenticate the separately requested DSS 15.0.2
server behavior or a complete supported client/server matrix.
\citep{exec24_workflow_contracts_dataiku}

The distinction matters even in convenience methods. The inspected
trigger-wait method with \texttt{no\_fail=True} returns \texttt{None}
when a trigger is cancelled. The surrounding \texttt{run\_and\_wait}
method passes that option into the trigger wait, then constructs a waiter
without checking whether a run exists. The waiter's next loop dereferences
\texttt{scenario\_run.run}. Under the documented cancelled-trigger
condition, this composition therefore has an unguarded \texttt{None} path.
This is a conditional static finding; no native execution, deployment
prevalence or observed customer failure is claimed. A wrapper should handle
a cancelled-before-start event explicitly instead of promising that the
permissive flag suppresses every exception. It can do so without porting
the whole scenario engine.~\citep{exec24_workflow_contracts_dataiku}

Once a run exists, the client treats the presence of its result as the end
of running. It exposes SUCCESS, WARNING, FAILED and ABORTED outcomes.
The strict completion path accepts only SUCCESS; WARNING is not silently
equivalent to it. A permissive return still needs an explicit consumer
decision about outcome and artifact validity. Backoff begins at 200
milliseconds, grows by a factor of 1.1 and caps each sleep at 15 seconds.
This per-sleep ceiling does not bound the number or duration of requests.
The inspected DSS HTTP delegate does not pass a timeout argument. These
facts motivate a separate observation deadline and transport policy; they
do not establish that every customized client session behaves identically.
\citep{exec24_workflow_contracts_dataiku}

For an authored example, suppose four requested triggers produce one
cancellation before a run, one successful accepted result, one warning
result outside the acceptance contract, and one failed run. There are four
requests, three actual runs and one acceptable artifact. If $A$ is the
number accepted, $N$ the number requested and $R$ the number started,
\begin{equation}
  \widehat p_{\mathrm{request}}=A/N=1/4,
  \qquad \widehat p_{\mathrm{started}}=A/R=1/3.
  \label{eq:exec24-workflow-denominator}
\end{equation}
Both rates are legitimate when labelled, but only the first retains the
cancelled request. With stipulated costs of USD 2, 4, 3 and 3, total cost
is USD 12: USD 4 per started run and USD 12 per accepted artifact. These
invented amounts demonstrate a denominator effect, not vendor prices or
bank economics. Three stipulated backoff sleeps total $331/500$ seconds.
With three two-second requests, observation takes $3331/500$ seconds before
other overhead. A sleep schedule is not a whole-workflow time guarantee.
\citep{exec24_workflow_contracts_reference}

The same distinction protects statistical evidence. Public-data evaluations
that retain only successful runs describe performance conditional on
execution and acceptance. They cannot represent all requested cases without
accounting for exclusions. The proposed CLI/MCP result should expose request,
trigger and run identities, whether execution began, native outcome,
exception and acceptance decision. A caller can then require strict success
or inspect a warning while retaining the original denominator.

\subsection{SAS: an execution example is not a supported server schema}

The pinned SAS repository describes its contents as preview REST examples.
Its Job Execution text describes asynchronous execution of a job definition,
submission through \texttt{/jobExecution/jobs}, and polling the created
job's state link. It distinguishes completed from failed and directs the
consumer to an error member on failure. Artifacts occupy a provider-specific
results map; the log location is separate. Without a supplied expiration,
the example assigns cleanup to the client. A workflow-origin submission
can identify its submitting application. These details support an adapter
that follows resource links and preserves provider identity, execution
state, artifacts and lifecycle actions separately.
\citep{exec24_workflow_contracts_sas}

The page identifies document version 5 and an update date of 26 November
2019, not a SAS server release. A cutoff repository revision authenticates
the bytes without making the examples a complete supported API schema.
The separate Visual Investigator workflow text concerns BPMN case management
for investigation and threat analysis. It must not be substituted for the
analytics workflow service used by an earlier Python client. Supported
editions, complete state/error enumerations, prompt schemas and timezone
behavior still need normative technical originals.
\citep{exec24_workflow_contracts_sas}

For the proposal, a narrow worker can bind the scientific target and data
revision to the definition, submission, provider and returned artifact
before remote cleanup. Reimplementing the REST stack would not remove
the missing server evidence. Public examples justify prototype interfaces
and test cases; they cannot supply licensed-deployment conformance results.

\subsection{Domino: method and failure class govern retry behavior}

The earlier Domino HTTP manager supplies a retry policy without overriding
its allowed-method set. Package setup permits Requests at least 2.4.2 and
urllib3 at least 1.26.19 but below 3; the requirements file narrows Requests
to the 2.31 line. A selected historical profile of Requests 2.31.0 and
urllib3 1.26.20 fits these inspected ranges, including Requests' own range.
It is a source-review profile, not a solved environment or deployment lock.
A separately acquired cutoff development revision cannot replace this
declared-range distinction.~\citep{exec24_workflow_contracts_domino}

Requests retains the supplied Retry object and passes it to the connection
pool. In the selected urllib3 profile, default allowed methods are HEAD,
GET, PUT, DELETE, OPTIONS and TRACE. POST and PATCH are absent. The response
retry predicate checks the method before its status list; the read-error
path also applies a method gate. Thus a retryable status list alone does
not imply that a submission POST is retried on those statuses. The
connection-error branch differs: it does not apply that gate. The source
classifies connection failures as cases in which the server is thought not
to have received the request, whereas read errors can follow processing.
A blanket claim that POST can never be retried would also be too broad.
\citep{exec24_workflow_contracts_domino}

A transport retry, authentication refresh, restart of a known remote job
and new scientific attempt must therefore remain distinct. After an
ambiguous submission, the proposed worker retains the unresolved remote
identity and reconciles it under an actual server idempotency contract.
HTTP defaults and response codes do not prove exactly-once execution.
The pinned public \texttt{api-docs} repository does not close this gap:
its original README and index redirect to a mutable latest guide rather
than storing a versioned \texttt{v4/jobs} schema. Complete terminal/error
semantics and the meaning of every conflict response remain unresolved.
\citep{exec24_workflow_contracts_domino}

The open client remains a plausible integration boundary because it already
exposes identities and job actions. A wholesale port would create a second
implementation with the same unresolved server assumptions. A prototype
should first lock full dependencies and register response-status,
connection, read, credential and reconciliation cases. No native trial or
claim that an installation duplicates or loses requests is made here.

\subsection{Posit: preserve a concrete backend's unsupported actions}

The C++ SDK's sample Local backend supplies a concrete execution contract.
It launches an asynchronous child, stores its process identifier and
publishes Pending. On exit it records the process exit code, then changes
every non-killed job to Finished without requiring zero. A quickly exiting
Pending job can first receive a Running notification. This backend explicitly
does not support cancellation. Kill, stop, suspend and resume use separate
signals and state paths. A proposed adapter must report unsupported Cancel
explicitly instead of turning it into a successful or merely skipped test.
\citep{exec24_workflow_contracts_posit}

Two stipulated records can both say Finished while one has exit zero and
a valid artifact and the other has exit seven and no acceptable artifact.
Acceptance needs the appropriate native state, an allowed exit code where
the contract uses one, a request-bound artifact and scientific validation.
A valid refusal may itself be an acceptable scientific answer if the contract
declares that possibility. The example describes invented records, not
observed bank jobs.~\citep{exec24_workflow_contracts_reference}

Historical documentation limits remain material. The prerequisite guide
mentions Workbench 1.4, later says 1.4 or greater, and requires a launcher-enabled
license and test scheduler. That statement is not a current support matrix.
The build specifies C++11, default Boost 1.83.0 and OpenSSL 1.1.1, which need
reconciliation with the older prerequisite table before building an environment.
The Local three-second setting concerns inter-node connections, not a global
causal-job deadline. The SDK notice distinguishes public permissive terms
from software received under a commercial agreement; it does not license
Workbench itself.~\citep{exec24_workflow_contracts_posit}

The smoke-test guide explicitly says its utility does not replace integration
testing with Launcher and the consuming product. It varies privilege and
configuration and requires manual verification of printed responses. Its
informal description of Finished as without failure must not erase the
executable exit-code distinction. A future profile should name exact product
and Launcher builds, SDK revision, backend, privilege mode, network/output
switches and allowed exit codes. Unsupported Cancel, races, missing artifacts
and failures each retain their own accounting. This is a testable prototype
plan, not completed conformance.~\citep{exec24_workflow_contracts_posit}

These comparisons favor small typed workers around existing execution
contracts. The causal library owns the intervention target, assumptions,
observation clock, artifact validation and failure accounting; platforms
retain their native execution vocabulary. Public evidence now supports a
more precise trigger/run distinction, retry classification and exit-code
contract. Historical server manuals, complete schemas, supported product
profiles, rights and measured costs still need their stated evidence. These
limits guide the next prototype without claiming complete commercial-workflow
coverage.
